\subsection{Mutual Information for Continuous Random Variables}

In the discrete setting, entropy is maximized by the uniform distribution: as the number of possible outcomes grows, so does the discrete entropy.  Extending this notion to continuous variables requires care, so we define mutual information via the Kullback–Leibler divergence:

\[
I(X;Y)
=\mathbb{E}\!\biggl[\log\frac{f_{X,Y}(X,Y)}{\,f_X(X)\,f_Y(Y)\,}\biggr]
= D\bigl(f_{X,Y}\,\|\,f_X\,f_Y\bigr).
\]

Equivalently, if we discretize $X$ and $Y$ into bins of width $\Delta$, obtaining $X^\Delta,Y^\Delta$, then
\[
I(X^\Delta;Y^\Delta)
=\mathbb{E}\!\Bigl[\log\frac{\Delta^2\,f_{X,Y}(X,Y)}{\Delta\,f_X(X)\,\Delta\,f_Y(Y)}\Bigr]
=\mathbb{E}\!\Bigl[\log\frac{f_{X,Y}(X,Y)}{f_X(X)\,f_Y(Y)}\Bigr],
\]
and letting $\Delta\to0$ recovers the continuous definition.

\subsection{Basic Properties of Continuous Mutual Information}

\begin{itemize}
  
  \item \textbf{Nonnegativity:}  
    \[    I(X;Y)\;=\;D\bigl(f_{X,Y}\,\|\,f_Xf_Y\bigr)\;\ge\;0,
    \]
    with equality iff $X$ and $Y$ are independent.
  \item \textbf{Symmetry:}  $I(X;Y)=I(Y;X)$.
  \item \textbf{Data processing inequality:}  
    If $X\to Y\to Z$ is a Markov chain, then
    \[
      I(X;Z)\le I(X;Y).
    \]
\end{itemize}




\subsection{Differential Entropy}

The \emph{differential entropy} of a continuous random variable $X$ with density $f_X(x)$ is
\[
h(X) = -\int f_X(x)\,\log f_X(x)\,\mathrm{d}x.
\]
If we discretize $X$ into bins of width $\Delta$, its discrete entropy satisfies
\[
H(X^\Delta)
=-\sum_i P\bigl(X\in[i,i+\Delta)\bigr)\log P\bigl(X\in[i,i+\Delta)\bigr)
= h(X) + \log\frac1\Delta + o(1),
\]
so $H(X^\Delta)\to\infty$ as $\Delta\to0$.

Unlike discrete entropy, differential entropy is not invariant under one-to-one maps.  In particular, for a nonzero scalar $a$,
\[
h(aX)=h(X)+\log|a|.
\]

It can even be negative. For example, the uniform distribution $U(0,1/2)$ has negative differential entropy;

$$\int_{0}^{1/2}-2\log2\,dx=-\log2<0.$$
Thus, differential entropy does not share desirable properties of discrete entropy. So, we generally use KL divergence or mutual information instead.

\subsection{Joint and Conditional Differential Entropy}

For a pair $(X_1,X_2)$ with joint density $f_{X_1,X_2}$, define
\[
h(X_1,X_2)
=-\iint f_{X_1,X_2}(x_1,x_2)\,\log f_{X_1,X_2}(x_1,x_2)\,\mathrm{d}x_1\mathrm{d}x_2,
\]
and the conditional entropy
\[
h(X_1\mid X_2)
=-\iint f_{X_1,X_2}(x_1,x_2)\,\log f_{X_1\mid X_2}(x_1\mid x_2)\,\mathrm{d}x_1\mathrm{d}x_2.
\]
These satisfy the chain rule:
\[
h(X_1,X_2)=h(X_2)+h(X_1\mid X_2).
\]

\subsubsection*{Invariance of mutual information}
Although $h(X)$ shifts under scaling, mutual information remains unchanged:
    \[
    I(aX;Y)
    = h(aX)-h(aX\mid Y)
    = \bigl(h(X)+\log|a|\bigr)-\bigl(h(X\mid Y)+\log|a|\bigr)
    = I(X;Y).
    \]

\subsection{Maximum Differential Entropy Under a Variance Constraint}

Just as the uniform distribution maximizes discrete entropy on a finite set, the Gaussian maximizes differential entropy among all densities with a given second moment.

\begin{theorem}
  If $\mathbb{E}[X^2]\le P$, then the Gaussian $\mathcal{N}(0,P)$ achieves the largest $h(X)$.
\end{theorem}

\begin{proof}
Let \(g(x)\) be a Gaussian PDF with mean \(\mu\) and variance \(P\) and
\(f(x)\) an arbitrary PDF with the same variance.  
(Since differential entropy is translation-invariant, we may assume that
\(f(x)\) has the same mean \(\mu\) as \(g(x)\).)

\medskip
\noindent
Consider the KL divergence between the two distributions:
\[
0 \;\le\; D_{KL}(f\|g)
      = \int_{-\infty}^{\infty} f(x)\,
        \log\!\Bigl(\frac{f(x)}{g(x)}\Bigr)\,dx
      = -h(f)-\int_{-\infty}^{\infty} f(x)\,\log(g(x))\,dx.
\]

\medskip
\noindent
Now note that
\begin{align*}
\int_{-\infty}^{\infty} f(x)\,\log(g(x))\,dx
&=\int_{-\infty}^{\infty}
  f(x)\,\log\!\Bigl(
    \frac{1}{\sqrt{2\pi P}}\,
    e^{-\frac{(x-\mu)^2}{2P}}
  \Bigr)dx \\[4pt]
&=\int_{-\infty}^{\infty}
  f(x)\,\log\frac{1}{\sqrt{2\pi P}}\,dx
  \;+\;\log(e)\int_{-\infty}^{\infty}
  f(x)\Bigl(-\frac{(x-\mu)^2}{2P}\Bigr)dx \\[4pt]
&=-\frac12\log(2\pi P)\;-\;\log(e)\,
      \frac{P}{2P} \\[4pt]
&=-\frac12\log(2\pi e P) \;=\; -\,h(g),
\end{align*}
because the result depends on \(f\) only through the common variance \(P\).
\medskip
\noindent
Combining the two results yields
\[
h(g)-h(f)\;\ge\;0,
\]
with equality when \(f(x)=g(x)\).
\end{proof}